\section*{Appendix a: Additional definitions of test fairness}
This appendix provides some details of fairness definitions included in Table~\ref{tab:fairness_criteria_summary} that were not introduced in the text 
of Section~\ref{sec:history}.

\subsection*{Einhorn and Bass}

In  1971, Einhorn and Bass 
\cite{einhorn1971methodological} noted that even if Cleary's criterion is satisfied, different rates of false positives and false negatives may be achieved for different subgroups due to differences in standard errors of estimate for the two subgroups.
That is, differences in variability around the common line of regression leads to different false positive and false negative rates. To address this, they propose a criterion based on achieving  equal false discovery rate, or as they put it, ``designated risk'', at the decision boundary. That is, $Prob(Y>y* | R=r_a*, A=a)$ is constant for all subgroups $a$.

\subsection*{Darlington's ``culturally optimum''}

Darlington \cite{darlington1971another} proposes that the
subjective value that one places on test validity (related to accuracy) and diversity can be scenario-specific.
He proposes a technique for eliciting these  value judgements, leading to a variable $k$ which measures the amount of tradeoff
in validity that is acceptable to increase diversity. He proposes
that the ``culturally optimum'' test is one that maximizes 
$\rho_{X(Y-kC)}$.

\subsection*{Jones}

In 1973, Jones  \cite{jones1973moderated}
proposed a ``general standard'' of fairness that is related to Thorndike's (and hence also related quota-based
definitions of fairness). In Jones' criterion, candidates are ranked in descending order both by test score and by ground truth. If an equal proportion of candidates from the subgroup are present in the top $n\%$ of both ranked lists then the test is fair ``at position $n$''. 
Jones' ``general standard'' of fairness requires that this hold for all values of $n$. Jones assumes a regression model relating test scores to ground truth, and also defines a weaker ``mean-fair'' criterion for a subgroup that
``the group's average predicted score
equals its average performance score on the [ground truth].''
